% GENERATED FILE. Edit canonical lessons, then run npm run generate:books.
% canonical-source-hash: fnv1a64:11fbba8639c9fabf
% canonical-lessons: ES-C649-moto, ES-C649-motorista, ES-C649-patinete, ES-C649-pedalear, ES-C649-ciclista

\chapter{Sobre dos ruedas}
\label{ch:es-sobre-dos-ruedas}
\hlchaptermodality{649}

\begin{chapteropening}
\textbf{By the end of this chapter:} \emph{I can talk about getting around on two wheels: motorbikes and their riders, scooters and bicycles.}

Complete the last lesson of chapter 649: I can talk about getting around on two wheels: motorbikes and their riders, scooters and bicycles.
\end{chapteropening}

\section[la moto]{la moto — the motorbike}

\label{lesson:ES-C649-moto}



\begin{quote}
Before the new one: say the Spanish for the lane, then the Spanish for to tow.
\end{quote}

\subsection*{You'll want to know: la moto}
\textbf{la moto} — \textquotedblleft{}the motorbike\textquotedblright{}.

\begin{quote}
Voy al trabajo en moto.
\end{quote}

\begin{quote}
\emph{I go to work by motorbike.}
\end{quote}

Short for \emph{la motocicleta}, so it is feminine.

\begin{cousinweb}
\emph{moto} is short for \emph{motocicleta}, which comes from French \emph{motocyclette} (EtymDB), built on words for an engine and a cycle.
\end{cousinweb}

\begin{grammarlens}[title={What you've built}]
One more word for telling, explaining and describing at B1.
\end{grammarlens}

\subsection*{Your turn}
\begin{itemize}
\raggedright
  \item \emph{Say it:} \emph{la moto}
  \item \emph{Say it:} \emph{la moto}, once more
  \item \emph{Say it:} a sentence of your own with \emph{la moto}
  \item \emph{Recall:} say the Spanish for the lane, then the Spanish for to tow, then say \emph{la moto} again
\end{itemize}

\begin{tcolorbox}[breakable,colback=teal!4,colframe=teal!35!black,title={Before you move on}]
What does la moto mean? (\textquotedblleft{}The motorbike\textquotedblright{}.) Say it once more.
\end{tcolorbox}
\section[el motorista]{el motorista — the motorcyclist}

\label{lesson:ES-C649-motorista}



\begin{quote}
Before the new one: say the Spanish for to tow, then the Spanish for the motorbike.
\end{quote}

\subsection*{You'll want to know: el motorista}
\textbf{el motorista} — \textquotedblleft{}the motorcyclist\textquotedblright{}.

\begin{quote}
Los motoristas llevan chaqueta y guantes.
\end{quote}

\begin{quote}
\emph{Motorcyclists wear jackets and gloves.}
\end{quote}

One form for both: \emph{el motorista}, \emph{la motorista}.

\begin{cousinweb}
\emph{motorista} is built on \emph{motor}, a learned borrowing from Latin \emph{mōtor} \textquotedblleft{}a mover\textquotedblright{} (Lewis \& Short), with \emph{-ista}, the ending for a person, from Greek \emph{-ist\'{\={e}}s}.
\end{cousinweb}

\begin{grammarlens}[title={What you've built}]
One more word for telling, explaining and describing at B1.
\end{grammarlens}

\subsection*{Your turn}
\begin{itemize}
\raggedright
  \item \emph{Say it:} \emph{el motorista}
  \item \emph{Say it:} \emph{el motorista}, once more
  \item \emph{Say it:} a sentence of your own with \emph{el motorista}
  \item \emph{Recall:} say the Spanish for to tow, then the Spanish for the motorbike, then say \emph{el motorista} again
\end{itemize}

\begin{tcolorbox}[breakable,colback=teal!4,colframe=teal!35!black,title={Before you move on}]
What does el motorista mean? (\textquotedblleft{}The motorcyclist\textquotedblright{}.) Say it once more.
\end{tcolorbox}
\section[el patinete]{el patinete — the scooter (kick or electric)}

\label{lesson:ES-C649-patinete}



\begin{quote}
Before the new one: say the Spanish for the motorbike, then the Spanish for the motorcyclist.
\end{quote}

\subsection*{You'll want to know: el patinete}
\textbf{el patinete} — \textquotedblleft{}the scooter (kick or electric)\textquotedblright{}.

\begin{quote}
Muchos jóvenes usan patinete eléctrico.
\end{quote}

\begin{quote}
\emph{A lot of young people use electric scooters.}
\end{quote}

In Latin America people often say \emph{el scooter} or \emph{el monopatín}.

\begin{cousinweb}
\emph{patinete} is built on \emph{patín} \textquotedblleft{}skate\textquotedblright{}, which was borrowed from French \emph{patin}, built on \emph{patte} \textquotedblleft{}paw, foot\textquotedblright{} (EtymDB).
\end{cousinweb}

\begin{grammarlens}[title={What you've built}]
One more word for telling, explaining and describing at B1.
\end{grammarlens}

\subsection*{Your turn}
\begin{itemize}
\raggedright
  \item \emph{Say it:} \emph{el patinete}
  \item \emph{Say it:} \emph{el patinete}, once more
  \item \emph{Say it:} a sentence of your own with \emph{el patinete}
  \item \emph{Recall:} say the Spanish for the motorbike, then the Spanish for the motorcyclist, then say \emph{el patinete} again
\end{itemize}

\begin{tcolorbox}[breakable,colback=teal!4,colframe=teal!35!black,title={Before you move on}]
What does el patinete mean? (\textquotedblleft{}The scooter (kick or electric)\textquotedblright{}.) Say it once more.
\end{tcolorbox}
\section[pedalear]{pedalear — to pedal}

\label{lesson:ES-C649-pedalear}



\begin{quote}
Before the new one: say the Spanish for the motorcyclist, then the Spanish for the scooter.
\end{quote}

\subsection*{You'll want to know: pedalear}
\textbf{pedalear} — \textquotedblleft{}to pedal\textquotedblright{}.

\begin{quote}
Pedaleo una hora cada mañana.
\end{quote}

\begin{quote}
\emph{I pedal for an hour every morning.}
\end{quote}

Built on \emph{el pedal}.

\begin{cousinweb}
\emph{pedalear} is built on \emph{pedal}, a learned borrowing from Latin \emph{pedālis} \textquotedblleft{}of the foot\textquotedblright{} (EtymDB), from \emph{pēs, pedis} \textquotedblleft{}foot\textquotedblright{}.
\end{cousinweb}

\begin{grammarlens}[title={What you've built}]
One more word for telling, explaining and describing at B1.
\end{grammarlens}

\subsection*{Your turn}
\begin{itemize}
\raggedright
  \item \emph{Say it:} \emph{pedalear}
  \item \emph{Say it:} \emph{pedalear}, once more
  \item \emph{Say it:} a sentence of your own with \emph{pedalear}
  \item \emph{Recall:} say the Spanish for the motorcyclist, then the Spanish for the scooter, then say \emph{pedalear} again
\end{itemize}

\begin{tcolorbox}[breakable,colback=teal!4,colframe=teal!35!black,title={Before you move on}]
What does pedalear mean? (\textquotedblleft{}To pedal\textquotedblright{}.) Say it once more.
\end{tcolorbox}
\section[el ciclista]{el ciclista — the cyclist}

\label{lesson:ES-C649-ciclista}



\begin{quote}
Before the new one: say the Spanish for the scooter, then the Spanish for to pedal.
\end{quote}

\subsection*{You'll want to know: el ciclista}
\textbf{el ciclista} — \textquotedblleft{}the cyclist\textquotedblright{}.

\begin{quote}
Hay muchos ciclistas en la carretera.
\end{quote}

\begin{quote}
\emph{There are a lot of cyclists on the road.}
\end{quote}

One form for both: \emph{el ciclista}, \emph{la ciclista}.

\begin{cousinweb}
\emph{ciclista} comes from French \emph{cycliste} (EtymDB), built on \emph{cycle}, from Greek \emph{kýklos} \textquotedblleft{}circle, wheel\textquotedblright{}.
\end{cousinweb}

\begin{grammarlens}[title={What you've built}]
One more word for telling, explaining and describing at B1.
\end{grammarlens}

\subsection*{Your turn}
\begin{itemize}
\raggedright
  \item \emph{Say it:} \emph{el ciclista}
  \item \emph{Say it:} \emph{el ciclista}, once more
  \item \emph{Say it:} a sentence of your own with \emph{el ciclista}
  \item \emph{Recall:} say the Spanish for the scooter, then the Spanish for to pedal, then say \emph{el ciclista} again
\end{itemize}

\begin{tcolorbox}[breakable,colback=teal!4,colframe=teal!35!black,title={Before you move on}]
What does el ciclista mean? (\textquotedblleft{}The cyclist\textquotedblright{}.) Say it once more.
\end{tcolorbox}
